% BHU PYQ -- Professional Exam Template
\documentclass[11pt,a4paper]{article}

\usepackage[a4paper,top=2.5cm,bottom=2.5cm,left=2.8cm,right=2.8cm,headheight=14pt]{geometry}
\usepackage{amsmath,amssymb}
\usepackage[T1]{fontenc}
\usepackage[utf8]{inputenc}
\usepackage{lmodern}
\usepackage{microtype}
\usepackage{fancyhdr}
\usepackage{enumitem}
\usepackage{listings}
\usepackage{hyperref}
\usepackage{lastpage}
\usepackage{parskip}

\hypersetup{colorlinks=true,urlcolor=black,linkcolor=black,
  pdftitle={BVCA-201: Introduction To Databases -- BHU PYQ}}

\pagestyle{fancy}
\fancyhf{}
\renewcommand{\headrulewidth}{0.4pt}
\renewcommand{\footrulewidth}{0.4pt}
\fancyhead[L]{\small\textit{Banaras Hindu University}}
\fancyhead[R]{\small\textit{BVCA-201: Introduction To Databases}}
\fancyfoot[C]{\small Page \thepage\ of \pageref{LastPage}}

\lstset{
  basicstyle=\small\ttfamily,
  frame=single,
  breaklines=true,
  columns=flexible,
  keepspaces=true
}

\newlist{parts}{enumerate}{2}
\setlist[parts,1]{label=(\alph*),leftmargin=*,itemsep=4pt,topsep=3pt}
\setlist[parts,2]{label=(\roman*),leftmargin=*,itemsep=2pt,topsep=2pt}
\newcommand{\pts}[1]{\hfill\textit{[#1]}}

\begin{document}
\begin{center}
  {\large\bfseries BANARAS HINDU UNIVERSITY}\\[2pt]
  {\normalsize B.Voc. Semester IV Examination, 2023-24}\\[6pt]
  \rule{0.8\linewidth}{0.5pt}\\[6pt]
  {\large\bfseries Paper: BVCA-201 --- Introduction To Databases}\\[4pt]
  {\normalsize Time: Three Hours \hfill Maximum Marks: 70}
\end{center}

\rule{\linewidth}{0.4pt}

\smallskip
\noindent\textit{Instructions: Please attempt all Sections. Question of Section — A, B and C carry 12, 6 and 1}

\bigskip

Reh-No. 76932 ,
: TT |
Subject: Computer Application
Note: Please attempt all Sections. Question of Section — A, B and C carry 12, 6 and 1
marks each respectively
SECTION --- A
Note: Please answer in approximately 500 words each. Attempt any two from the following
question. Each question carries 12 Marks
‘ Question 1. Describe the three-schema architecture. Why do we need mappings between
schema levels? How do different schema definition languages support this architecture?

\medskip\noindent\textbf{Question 2.} Discuss the different types of user-friendly interfaces and the types of users who
typically use each.

\medskip\noindent\textbf{Question 3.} Discuss the various reasons that lead to the occurrence of NULL values in
relations

\medskip\noindent\textbf{Question 4.} Discuss the role of a high-level data model in the database design process.

\bigskip\noindent{\large\bfseries SECTION --- B}
\rule{\linewidth}{0.4pt}\smallskip

Note: Please answer approximately in 200 words each. Each question carries 6 marks each.
| Attempt any six from following:
| Question 5.
A. What is the difference between logical data independence and physical data
independence? Which one is harder to achieve? Why?
B. Discuss the different data types of SQL.
C. Discuss the entity integrity and referential integrity constraints. Why each is considered
important?
D. Discuss Functional dependency with suitable example.
E. What is a relationship type? Explain the differences among a relationship instance, a
relationship type, and a relationship set.
F, What is a participation role? When is it necessary to use role names in the description of
relationship types?
G. What is meant by a recursive relationship type? Give some examples of recursive
relationship types.
H. Why are duplicate tuples not allowed in a relation?
I. What is the difference between a key and a superkey?
EEG
)

\bigskip\noindent{\large\bfseries SECTION --- C}
\rule{\linewidth}{0.4pt}\smallskip

Note: Please answer in briefest possible way. Each question carries 1 mark each.

\medskip\noindent\textbf{Question 6.} Define the following terms
\begin{parts}
  \item Data model
  \item Database schema
  \item Client/server architecture
  \item Owner entity type :
  \item Weak entity type i
  \item Entity
g. Attribute Value
  \item Relation state
  | i. Degree of a relation .
  \item Normalization
  —\_—\_— Oo -
\end{parts}

\vfill
\rule{\linewidth}{0.4pt}
\begin{center}\small
\href{https://bhu-exam.vercel.app/}{Click here to visit our website}\\[4pt]\href{https://me-aryan.vercel.app/}{Click here to visit our contributor}\\[4pt]\href{https://quashan-me.vercel.app/}{Click here to visit our contributor}
\end{center}
\end{document}
